% ===================================================================
%  GAMBAR No. 6 — Jero Omah, Perabot, Pangan, Padhang.
%
%  Traduction, faite en 2026, en javanais d'aujourd'hui, sur l'ido de
%  texto/io. Regles de la colonne : voir 10-gambar-01.tex. Cinq
%  numerotations, une par piece : t06-c1-* la chambre, c2 le salon,
%  c3 la salle a manger, c4 la cuisine, c5 la salle de bains.
%
%  RETOUR AU RECIT, ET LA REGLE DE NIVEAU REDEVIENT ACTIVE. Le tableau
%  5 etait un dialogue et kolonoj.py y suspendait les huit regles de
%  krama ; celui-ci ne porte aucune attribution de parole, le compte
%  retombe a zero, et « narracio » s'applique de nouveau sans qu'on ait
%  rien a dire. Le seuil pose a cinq n'a pas eu a etre retouche : c'est
%  la premiere verification en vraie grandeur du mecanisme.
%
%  LE PLUS GRAND NOMBRE DE RENVOIS DU LIVRET — DEUX CENT ONZE — ET
%  PRESQUE QUE DES OBJETS. C'est le tableau du mobilier, et il partage
%  le vocabulaire de la maison en deux moities tres nettes.
%
%   * CE QU'ON POSE, CE QU'ON FAIT CUIRE, CE QU'ON BALAIE EST JAVANAIS :
%     jogan le plancher, pyan le plafond, kasur le matelas (22), kemul
%     la couverture (24), bantal l'oreiller (26), klasa la natte, sapu
%     le balai (34), kendhi la cruche (50), wajan la poele (47), kendhil
%     la marmite (46), kukusan, irus la louche (21), enthong, talenan la
%     planche a hacher (41), gandhik, pawon le foyer (1), geni le feu
%     (3), areng le charbon, sumur, jedhing.
%   * CE QUI EST ARRIVE PAR BATEAU EST NEERLANDAIS : kasten l'armoire,
%     spiegel le miroir (60), ember le seau (9), sprei le drap (23),
%     lampu, kompor le fourneau (19), oven le four (20), teko la
%     theiere (31), kraan le robinet (48), wastafel, dhus la douche
%     (7), gas, korek l'allumette (12).
%
%  ET LE PARTAGE N'EST PAS UNE QUESTION DE PRESTIGE, C'EST UNE QUESTION
%  DE DATE : tout ce qu'une maison javanaise avait deja porte son nom
%  javanais, tout ce que le Hollandais a installe porte le sien. La
%  cuisine en est la preuve, parce qu'elle contient les deux a la fois
%  — le pawon de terre et le kompor a gaz sont sur la meme planche, et
%  chacun garde le mot de son siecle.
%
%  wajan (47) EST LE SEUL MOT DE CE TABLEAU QUI SOIT SORTI DE JAVA.
%  La poele javanaise a donne son nom au wok chinois dans les langues
%  d'Europe par le malais et le neerlandais, et l'anglais dit « wok »
%  mais le neerlandais dit « wadjan ». C'est le troisieme mot du releve
%  qui voyage dans ce sens-la, apres قلی et بازار de la colonne
%  ourdoue : presque tout ce livret decrit des choses qui sont venues,
%  et il faut un tableau de cuisine pour qu'une chose parte.
%
%  LE JAVANAIS A UN MOT POUR LE BERCEAU (61) QUE NI L'IDO NI LE
%  FRANCAIS N'ONT : « kloso ayun » ne le rend pas, c'est « gendhongan »
%  qui le rend — le berceau suspendu qu'on balance d'une main en
%  travaillant de l'autre. Mais la planche grave un berceau europeen a
%  bascule, sur pieds, et la colonne ecrit donc « ayunan bayi », qui
%  decrit ce qui est grave. Meme refus qu'au gobak sodor du tableau 2
%  et au joged du tableau 3 : le mot d'ici ne se pose pas sur une chose
%  de la-bas, meme quand il est meilleur.
% ===================================================================

%%K t06-apar-1 apar
\VUsaut{6.0mm}
\VUtitre{15.0pt}{60}{GAMBAR No. 6}
\VUsaut{1.4mm}
\VUtitre{11.4pt}{0}{Jero Omah, Perabot, Pangan,}
\VUsaut{0.4mm}
\VUtitre{11.4pt}{0}{Padhang.}
\VUsaut{2.2mm}

%%K t06-apar-2 apar
Omah iku wis rampung. Pakdhene Ioannes wis manggon ing omah anyare,
kang tata ruwange wis kita ngerteni. Saiki ayo padha ndeleng kepriye
dheweke ngiseni omahe nganggo perabot. Kang bakal kita jlajahi
dhisik :

%%K t06-tit-1 sub
\VUpk{10.2pt}{40}{Kamar Turu.}
\VUsaut{3.0mm}

%%K t06-c1-01-1 p
1. --- Kita teka ing wektu kang kurang trep, amarga batur wadon lagi
sibuk ngresiki \VUgras{kamar} iku.

%%K t06-c1-02-1 p
2. --- Ing \VUgras{meja rias} \textsuperscript{(1)} ana ing kene lan kana
warna-warna barang kanggo wong wisuh, nyisir, cekake kanggo dandan.
Yaiku \VUgras{baskom} \textsuperscript{(2)} lan \VUgras{kendhi banyu}
\textsuperscript{(3)}, \VUgras{sikat rambut} \textsuperscript{(4)},
\VUgras{jungkat} \textsuperscript{(5)}, \VUgras{sabun}
\textsuperscript{(6)} lan \VUgras{spons} \textsuperscript{(7)}, sarta
\VUgras{botol cilik} \textsuperscript{(8)} kanggo obat kumur untu.

%%K t06-c1-03-1 p
3. --- Ing jogan, ing ngarepe meja rias, iki \VUgras{ember}
\textsuperscript{(9)} kanggo nampani banyu reged.

%%K t06-c1-04-1 p
4. --- Ing \VUgras{ruwang ngarep kamar} \textsuperscript{(10)} kita weruh
sawetara \VUgras{cantholan klambi} \textsuperscript{(13)} lan
\VUgras{payung} \textsuperscript{(11)} loro ing \VUgras{wadhah payung}
\textsuperscript{(12)}.

%%K t06-c1-05-1 p
5. --- Ing sisih sijine lawang ana \VUgras{lemari pangilon}
\textsuperscript{(14)} kang isi \VUgras{sandhangan lawon}
\textsuperscript{(15)}.

%%K t06-c1-06-1 p
6. --- Ayo kita nganggo kesempatan nalika \VUgras{batur wadon}
\textsuperscript{(6)} nata \VUgras{peturon} \textsuperscript{(17)}, kanggo
njlajahi perangan-perangane. \VUgras{Rangka peturon}
\textsuperscript{(20)} isi \VUgras{kasur per} \textsuperscript{(21)} kang
becik, lan ing dhuwure Iustina sedhela maneh bakal nyelehake
\VUgras{kasur} \textsuperscript{(22)} wol. Banjur dheweke bakal njupuk
saka kursi \VUgras{sprei} \textsuperscript{(23)} loro lan \VUgras{kemul}
\textsuperscript{(24)}. Sawise iku dheweke bakal masang ing sirahe
peturon \VUgras{guling melintang} \textsuperscript{(25)} lan
\VUgras{bantal} \textsuperscript{(26)}, kang saiki isih ana ing dhuwur
\VUgras{permadani peturon} \textsuperscript{(32)} bareng karo
\VUgras{kemul wulu} \textsuperscript{(27)}. Dheweke nganggo sulak wulu
kanggo ngresiki \VUgras{gordhen} \textsuperscript{(19)} lan
\VUgras{kelambu peturon} \textsuperscript{(18)}, lan iki peranganing
pagawean wis rampung.

%%K t06-c1-07-1 p
7. --- Banjur dheweke kudu rada nata \VUgras{meja cilik peturon}
\textsuperscript{(28)}, muteri \VUgras{weker} \textsuperscript{(29)},
nyiyapake \VUgras{lampu wengi} \textsuperscript{(30)}, njupuk
\VUgras{lilin} \textsuperscript{(31)} kanggo ngresiki cagak lilin.

%%K t06-tit-2 sub
\VUsaut{5.0mm}
\VUpk{10.2pt}{40}{Kranjang jaitan lan piranti pawon geni.}
\VUsaut{3.0mm}

%%K t06-c1-08-1 p
8. --- \VUgras{Kranjang jaitan} \textsuperscript{(34)} ing cedhake
\VUgras{dhingklik} \textsuperscript{(33)} uga butuh banget ditata.
Dheweke kudu nglempit \VUgras{bordhiran} \textsuperscript{(35)},
nglebokake ing \VUgras{meja jaitan} \textsuperscript{(38)}
\VUgras{tudhung driji,} \VUgras{benang} \textsuperscript{(36)},
\VUgras{dom} \textsuperscript{(37)} lan \VUgras{gunting}
\textsuperscript{(39)}, banjur ngunci tutupe meja iku nganggo
\VUgras{kunci} \textsuperscript{(40)}.

%%K t06-c1-09-1 p
9. --- Ing \VUgras{meja sikil siji} \textsuperscript{(41)} ing tengah,
Iustina bakal ngganti banyune \VUgras{karaf} \textsuperscript{(42)}.
Dheweke bakal nyelehake \VUgras{gula} ing \VUgras{wadhah gula}
\textsuperscript{(43)}, ngresiki \VUgras{sapit gula}
\textsuperscript{(44)} lan njupuk \VUgras{sikat sandhangan}
\textsuperscript{(46)}.

%%K t06-c1-10-1 p
10. --- Sisih tengene kamar bakal njaluk pagawean kang luwih sethithik
saka dheweke, amarga \VUgras{kursi dawa} \textsuperscript{(47)} lan
\VUgras{bantale} \textsuperscript{(48)} wis ana ing panggonan kang trep,
lan, sarehne hawane ora adhem, ora ana kang owah ing antarane piranti
\VUgras{pawon geni} \textsuperscript{(49)}. \VUgras{Sekop}
\textsuperscript{(50)}, \VUgras{sapit gedhe} \textsuperscript{(51)},
\VUgras{cagak kayu} \textsuperscript{(55)} lan \VUgras{pager awu}
\textsuperscript{(56)} wis resik ; \VUgras{tameng geni}
\textsuperscript{(54)} ora diowahi, lan \VUgras{kothak kayu}
\textsuperscript{(52)} wis kebak \VUgras{kayu obong}
\textsuperscript{(53)}.

%%K t06-noto-1 noto
(*) Babo = bocah cilik ; I. baby, P. bébé, It. bambino.

%%K t06-noto-2 noto
(*) Lambrequino = P. lambrequin, Sp. lambrequines, Port. lambrequins,
malah J. I. lambrequins ; dadi J. I. P. Port. Sp.

%%K t06-c1-11-1 p
11. --- Sanajan mangkono dheweke isih kudu ngresiki lebu \VUgras{jam
bandhul} \textsuperscript{(57)}, \VUgras{cagak lilin gedhe}
\textsuperscript{(58)} lan \VUgras{wadhah barang cilik}
\textsuperscript{(59)}, pungkasane ngelapi \VUgras{pangilon}
\textsuperscript{(60)}.

%%K t06-c1-12-1 p
12. --- Bab \VUgras{ayunan bayi} \textsuperscript{(61)} kanggo bayi cilik
— babo (*) —, iku wis siyap.

%%K t06-tit-3 sub
\VUsaut{5.0mm}
\VUpk{10.2pt}{40}{Ruwang Tamu.}
\VUsaut{3.0mm}

%%K t06-c2-01-1 p
1. --- Saiki iki ruwang tamu. Iku \VUgras{kamar} kang endah banget.
Pawon geni, kang ana ing pojok tengen, direngga \VUgras{reca cilik}
\textsuperscript{(1)} kang alus, \VUgras{jambangan}
\textsuperscript{(3)} endah loro, lan katutupan \VUgras{lambrekin}
\textsuperscript{(2)} beludru (*).

%%K t06-c2-02-1 p
2. --- Supaya kelethikan geni saka pawon ora ngobong permadani, ing
ngarepe pawon dipasang \VUgras{panulak kelethikan} \textsuperscript{(4)}
tembaga.

%%K t06-c2-03-1 p
3. --- Ing cedhake, \VUgras{meja tulis} \textsuperscript{(5)} wanita kang
apik lagi menga. Ing dhuwure \VUgras{nyonyah omah}
\textsuperscript{(23)} nulis layang-layange.

%%K t06-c2-04-1 p
4. --- Jendhelane direngga \VUgras{gordhen kaca} \textsuperscript{(6)}
karo \VUgras{taline} \textsuperscript{(8)} kang dicanthelake ing
\VUgras{canthelan} \textsuperscript{(7)}, lan gordhen gedhe saka lawon
kang \VUgras{dilempit ing dhuwur} \textsuperscript{(9)}.

%%K t06-c2-05-1 p
5. --- Ing cerukane jendhela, \VUgras{wadhah pot} \textsuperscript{(11)}
isi \VUgras{tanduran} \textsuperscript{(10)} ijo, yaiku palem kang endah
kang godhonge megar meh tekan \VUgras{lampu lenga}
\textsuperscript{(12)}.

%%K t06-c2-06-1 p
6. --- \VUgras{Tudhung lampu} \textsuperscript{(13)} lampu iku ditata
Maria, \VUgras{prawan} \textsuperscript{(14)} kang ayu, kang ana ing
\VUgras{piano} \textsuperscript{(15)}. Lungguh jejeg banget ing
\VUgras{dhingklik} \textsuperscript{(16)}, dheweke sinau lelagon. Dheweke
trampil lan ngati-ati banget, kaya kabukten saka \VUgras{lemari
musike} \textsuperscript{(17)} kang tinata rapi.

%%K t06-tit-4 sub
\VUsaut{5.0mm}
\VUpk{10.2pt}{40}{Bocah Nakal.}
\VUsaut{3.0mm}

%%K t06-c2-07-1 p
7. --- Adhine lanang, Paulus, nggoleki \VUgras{buku}
\textsuperscript{(19)} ing \VUgras{rak buku} \textsuperscript{(18)}.
Dheweke isih \VUgras{nom-noman} \textsuperscript{(20)}, ora bisa meneng
lan angel dieman. Karo gandhulan ing rak buku supaya bisa nggayuh buku
kang dhuwur banget, dheweke kentekan bebaya nibakake \VUgras{reca
dhadha} \textsuperscript{(21)} kang ana ing ndhuwure.

%%K t06-c2-08-1 p
8. --- Dheweke nganggo kesempatan nglakoni bab bodho iku nalika ibune
lagi repot ngobrol karo kancane, yaiku \VUgras{nyonyah}
\textsuperscript{(22)} kang teka nginep sawetara dina ing omahe. Ibune
lungguh ing \VUgras{kursi dawa} \textsuperscript{(24)} cedhak
\VUgras{kursi males} \textsuperscript{(25)}, dene kancane lungguh ing
\VUgras{kursi} \textsuperscript{(27)} kang mung katon
\VUgras{sandharane} \textsuperscript{(26)}.

%%K t06-c2-09-1 p
9. --- \VUgras{Pigura} \textsuperscript{(30)} gedhe kang gumantung ing
tembok isi \VUgras{potret} \textsuperscript{(29)} sanak wadon. Ing
\VUgras{album} \textsuperscript{(32)} kang diselehake ing meja
kaklumpukake \VUgras{potret-potret} \textsuperscript{(33)} sanak kadang
liyane. Bocah-bocah lagi wae mbukaki album iku, amarga
\VUgras{kursi-kursi} \textsuperscript{(31)} isih nglumpuk ing sakubenge
meja.

%%K t06-c2-10-1 p
10. --- \VUgras{Jambangan Cina} \textsuperscript{(34)} porselen, kang ana
ing cedhake \VUgras{ladhing kertas} \textsuperscript{(35)}, dianggo
kado marang Maria nalika pestane, lan \VUgras{kelir lempitan}
\textsuperscript{(36)} kang ngrengga pojoke kamar digawa saka Cina
dening pakdhene.

%%K t06-c2-11-1 p
11. --- Kabungahake dening \VUgras{lukisan} \textsuperscript{(37)} endah
kang gumantung ing \VUgras{tembok} \textsuperscript{(42)}, kapadhangan
dening \VUgras{lampu gantung pyan} \textsuperscript{(38)}, kang
\VUgras{lampu listrike} \textsuperscript{(39)} mencorong bungah ing
wayah bengi, ruwang tamu iku katon nyenengake banget.
\VUgras{Permadani} \textsuperscript{(40)} kang alus, kang kegelar ing
jogan papan, lan \VUgras{bel listrik} \textsuperscript{(41)} kang
gampang temen dianggo, nyampurnakake kepenake.

%%K t06-tit-5 sub
\VUsaut{3.6mm}
\VUpk{10.2pt}{40}{Ruwang Mangan.}
\VUsaut{3.0mm}

%%K t06-c3-01-1 p
1. --- Lonceng bujana bengi wis muni. Kabeh kulawarga, kejaba Maria,
wis nglumpuk ing sakubenge meja. Ruwang mangan iku anget kepenak
dening \VUgras{tungku} \textsuperscript{(1)} fayans kang becik banget,
karo \VUgras{cerobong} \textsuperscript{(2)} kang tipis.

%%K t06-c3-02-1 p
2. --- Kamar iku ora akeh perabote. Sadawane tembok gumantung, ing
dhuwure \VUgras{bangku cilik} \textsuperscript{(4)}, \VUgras{pancuran
cilik} \textsuperscript{(3)} kang mung kanggo rerenggan. Cedhak banget,
ana \VUgras{lemari pinggir} \textsuperscript{(5)}, panggonan wong
nyelehake panganan adhem, piring panganan panutup, lan warna-warna
piranti meja kang durung dibutuhake.

%%K t06-c3-03-1 p
3. --- Mangkono kita weruh ing rak ndhuwur \VUgras{pitik panggang}
\textsuperscript{(7)}, \VUgras{iwak} \textsuperscript{(8)} lan
\VUgras{mangkok} \textsuperscript{(10)} isi \VUgras{woh-wohan}
\textsuperscript{(9)}. Ing ngisore iki \VUgras{keju}
\textsuperscript{(11)}, \VUgras{roti} \textsuperscript{(12)}, lan
\VUgras{wadhah botol} \textsuperscript{(13)} kang isi kabeh kang
dibutuhake kanggo mbumboni salad : lenga, cuka, \VUgras{uyah}
\textsuperscript{(14)} lan \VUgras{mrica} \textsuperscript{(15)}.
Pungkasane, ing rak ngisor ana \VUgras{roti tar} \textsuperscript{(17)}
cedhak \VUgras{wadhah mosterd} \textsuperscript{(16)}.

%%K t06-c3-04-1 p
4. --- Jam ruwang mangan, kang diarani \VUgras{jam tembok}
\textsuperscript{(20)}, wangune aneh banget. Jam iku kepasang ing
pigura kayu kang uga isi \VUgras{barometer} \textsuperscript{(19)} lan
\VUgras{termometer} \textsuperscript{(18)}.

%%K t06-tit-6 sub
\VUsaut{4.6mm}
\VUpk{10.2pt}{40}{Ngladeni Bujana Bengi.}
\VUsaut{3.0mm}

%%K t06-c3-05-1 p
5. --- Ing kabeh temboke kamar dipasang \VUgras{papan kayu}
\textsuperscript{(21)} saka kayu jati kang dililini. \VUgras{Gordhen
lawang} \textsuperscript{(22)} saka lawon nutup cukup rapet lawang kang
tumuju emper. Ing kono kulawarga bakal ngombe kopi sawise bujana
bengi. \VUgras{Cangkir} \textsuperscript{(24)} lan \VUgras{lepek}
\textsuperscript{(25)} wis siyap ing \VUgras{lemari piring}
\textsuperscript{(23)}.

%%K t06-c3-06-1 p
6. --- Ing dhuwure meja kepasang ing pyan \VUgras{lampu gantung}
\textsuperscript{(26)}, kang lampune mencorong banget lan ing wayah
bengi madhangi saruwang mangan.

%%K t06-c3-07-1 p
7. --- Saiki ayo kita jlajahi \VUgras{meja mangan}
\textsuperscript{(27)} iku dhewe. Nalika kulawarga mlebu, piranti
mangan wis siyap. Ing \VUgras{taplak meja} \textsuperscript{(28)} wis
diselehake, ing panggonane saben wong kang bakal mangan,
\VUgras{piring} \textsuperscript{(33)}, \VUgras{serbet}
\textsuperscript{(29)}, \VUgras{sendhok} \textsuperscript{(34)},
\VUgras{porok} \textsuperscript{(35)}, \VUgras{ladhing}
\textsuperscript{(36)} lan \VUgras{gelas} \textsuperscript{(32)}. Ana uga
\VUgras{botol} \textsuperscript{(30)} kanggo anggur lan \VUgras{karaf}
\textsuperscript{(31)} kanggo banyu.

%%K t06-c3-08-1 p
8. --- \VUgras{Pelayan} \textsuperscript{(39)} kang kapisan nggawa
\VUgras{wadhah sup} \textsuperscript{(37)}, kang isih kumelun. Saiki
dheweke ngladekake \VUgras{panganan} \textsuperscript{(38)} kapisan,
yaiku panganan daging. Banjur dheweke bakal nyuguhake panganan kang wis
kita sebut, lan pungkasane, sawise panganan panutup, dheweke bakal
nganggo kesempatan nalika para bendarane wis mlebu emper, kanggo
ngangkati piranti meja lan nata maneh ruwang mangan.

%%K t06-c3-08-2 p
\textit{Cathetan.} --- \textit{Tembung lumrah kang gegayutan karo
pangan mung disebut cekak ing kene. Dhaptare bakal dijangkepi ing
gambar loro « Hotel » (No. 13) lan « Pasar » (No. 15).}

%%K t06-tit-7 sub
\VUsaut{4.6mm}
\VUpk{10.2pt}{40}{Pawon.}
\VUsaut{3.0mm}

%%K t06-c4-01-1 p
1. --- Ing pawon, kang kita deleng liwat lawang kang menga separo,
kita weruh rerusuhan kang endah. Ing ngarepe \VUgras{pawon geni}
\textsuperscript{(1)}, panggonan \VUgras{geni} \textsuperscript{(3)}
krepyak-krepyak, dipasang \VUgras{puteran sunduk}
\textsuperscript{(5)}. \VUgras{Panggangan} \textsuperscript{(7)} kang
endah wis \VUgras{disunduk} \textsuperscript{(6)} lan lagi mateng.

%%K t06-c4-02-1 p
2. --- \VUgras{Sebul geni} \textsuperscript{(8)} lan \VUgras{sekop
areng} \textsuperscript{(9)}, kang lagi wae dianggo, isih gumlethak ing
jogan bareng karo \VUgras{kayu obong} \textsuperscript{(10)}.

%%K t06-c4-03-1 p
3. --- Dhuwure pawon geni wis tinata : \VUgras{lentera}
\textsuperscript{(11)}, \VUgras{wadhah korek} \textsuperscript{(13)}
kang isi \VUgras{korek} \textsuperscript{(12)}, \VUgras{cagak lilin}
\textsuperscript{(14)} lan \VUgras{wadhah lilin} \textsuperscript{(15)}
karo \VUgras{liline} \textsuperscript{(16)} ana ing panggonane. Nanging
\VUgras{ember areng} \textsuperscript{(18)} kang gedhe, kebak
\VUgras{areng watu} \textsuperscript{(17)} kang lagi wae diisi
\VUgras{juru masak} \textsuperscript{(23)} ing kamar ngisor lemah, isih
kari ing tengahe pawon.

%%K t06-c4-04-1 p
4. --- Pancen, wong wadon kang mesakake iku kudu gage-gage supaya
sarapan siyap ing wektune. Saiki dheweke ana ing cedhake
\VUgras{kompor} \textsuperscript{(19)} lan ngudhek \VUgras{saus}
\textsuperscript{(24)} kang ora kena gosong banget.

%%K t06-c4-05-1 p
5. --- Ing \VUgras{oven} \textsuperscript{(20)} lagi mateng roti tar,
kang disiyapake kanggo panganan cilike bocah-bocah. Ing dhuwure kompor
gumantung \VUgras{irus} \textsuperscript{(21)} lan \VUgras{serok}
\textsuperscript{(22)}.

%%K t06-tit-8 sub
\VUsaut{4.0mm}
\VUpk{10.2pt}{40}{Piranti Pawon.}
\VUsaut{3.0mm}

%%K t06-c4-06-1 p
6. --- \VUgras{Piranti masak} \textsuperscript{(25)}, kang dicanthelake
ing pungkasane kamar, resik banget. \VUgras{Panci}
\textsuperscript{(26)} tembaga kang endah, \VUgras{poci susu}
\textsuperscript{(27)}, \VUgras{ayakan cilik} \textsuperscript{(28)} lan
\VUgras{parutan} \textsuperscript{(29)} wis diresiki kanthi ngati-ati.

%%K t06-c4-07-1 p
7. --- \VUgras{Kothak uyah} \textsuperscript{(30)} diselehake ing lemari
piring cedhak \VUgras{teko} \textsuperscript{(31)} lan \VUgras{botol
gedhe lenga} \textsuperscript{(32)}. \VUgras{Laci}
\textsuperscript{(33)} mesthine isi piranti mangan lan ladhing pawon.

%%K t06-c4-08-1 p
8. --- \VUgras{Sapu} \textsuperscript{(34)} lan \VUgras{sulak wulu}
\textsuperscript{(35)} ana ing pojok. Juru masak lagi wae nganggo
\VUgras{blok kayu} \textsuperscript{(36)}, banjur ditinggal ing cedhake
jendhela.

%%K t06-c4-09-1 p
9. --- \VUgras{Lap tangan} \textsuperscript{(37)} gumantung ing tembok,
cedhak \VUgras{corong} \textsuperscript{(38)}. Ing perangan pawon kono
disimpen \VUgras{panggangan kawat} \textsuperscript{(39)},
\VUgras{peso pencacah} \textsuperscript{(40)} lan \VUgras{talenan}
\textsuperscript{(41)}. Banjur, ing dhuwure peso pencacah, ing
\VUgras{rak} \textsuperscript{(42)}, ana \VUgras{kuwali}
\textsuperscript{(43)}, \VUgras{kendhil} \textsuperscript{(44)} karo
\VUgras{tutupe} \textsuperscript{(45)} lan \VUgras{dandang}
\textsuperscript{(46)}. \VUgras{Wajan} \textsuperscript{(47)} gumantung
ing cedhake.

%%K t06-c4-10-1 p
10. --- Mung \VUgras{kran} \textsuperscript{(48)} tembaga kang
mencorong ing pojok kene. \VUgras{Bak cuci} \textsuperscript{(49)},
panggonan \VUgras{kendhi gedhe} \textsuperscript{(50)} kang kandel
diselehake, mesthine kepenak banget karo lemari cilike, panggonan wong
nyimpen \VUgras{tong cuka} \textsuperscript{(51)} kang duwe
\VUgras{krane} \textsuperscript{(52)}, \VUgras{kothak sampah}
\textsuperscript{(53)} lan \VUgras{piring reged} \textsuperscript{(54)}.

%%K t06-tit-9 sub
\VUsaut{4.0mm}
\VUpk{10.2pt}{40}{Barang liya ing pawon.}
\VUsaut{3.0mm}

%%K t06-c4-11-1 p
11. --- \VUgras{Bak gedhe} \textsuperscript{(55)}, panggonan wong ngumbah
\VUgras{sayuran} \textsuperscript{(58)}, lan \VUgras{kwali wesi}
\textsuperscript{(56)} kang diselehake ing \VUgras{jogan ubin}
\textsuperscript{(57)} mesthine uga duwe panggonan ing lemari iku.

%%K t06-c4-12-1 p
12. --- Juru masak mesthi ora kekurangan tungku, amarga iki ana maneh,
ing pojoke meja, \VUgras{kompor gas} \textsuperscript{(59)} kang
dianggo manasi banyu ing panci cilik. \VUgras{Kukus}
\textsuperscript{(60)} kang metu saka kono nuduhake marang kita yen
banyune mesthi wis umob. Kaya-kaya banyu iku bakal dianggo kanggo kopi,
amarga iki ing pucuk sijine meja ana \VUgras{poci kopi}
\textsuperscript{(64)} lan \VUgras{gilingan kopi} \textsuperscript{(63)}.

%%K t06-c4-13-1 p
13. --- Kompor iku disambung karo \VUgras{kran gas}
\textsuperscript{(62)} nganggo \VUgras{selang karet}
\textsuperscript{(61)} kang dawa.

%%K t06-c4-14-1 p
14. --- \VUgras{Batur harian} \textsuperscript{(66)} kang teka mbantu
juru masak lagi nyiyapake sayuran kanggo bujana bengi. Yen sayuran iku
wis diresiki lan diumbah, dheweke bakal nyelehake ing \VUgras{baskom}
\textsuperscript{(65)} ngenteni wektune digodhog.

%%K t06-tit-10 sub
\VUsaut{4.0mm}
{\centering\fontsize{10.2pt}{10.2pt}\selectfont\textls[40]{\scshape Kamar Adus.} \textit{(Papan adus.)}\par}
\VUsaut{3.0mm}

%%K t06-c5-01-1 p
1. --- Kari siji bab kang rada ora sopan kanggo ngerteni kamar-kamar
baku ing omah iki : ndeleng tatanane kamar adus. Iku ora njaluk wektu
dawa, amarga kamare ora amba, lan kita enggal ngerti yen \VUgras{bak
adus} \textsuperscript{(1)} duwe \VUgras{pemanas banyu adus}
\textsuperscript{(2)} kang becik, yen \VUgras{wadhah sabun}
\textsuperscript{(3)} bisa digayuh tangane wong kang adus, lan yen
\VUgras{cantholan andhuk} \textsuperscript{(5)} mesthi nggampangake
garinge \VUgras{andhuk} \textsuperscript{(4)}.

%%K t06-c5-02-1 p
2. --- Kamar iku temboke katutupan \VUgras{ubin fayans}
\textsuperscript{(6)}, kang mbisakake masang \VUgras{piranti dhus}
\textsuperscript{(7)} tanpa wedi babar pisan yen banyu kang muncrat
nalika dianggo bakal ngrusak tembok. \VUgras{Baskom cilik}
\textsuperscript{(8)} seng, kang kanggo ngedusi sikil, njangkepi
perabot iku.

\VUsaut{6.0mm}
\VUfilet{22mm}
